\markboth{TREE}{TREE}\label{sec:06}
\textbf{Binary Lifting}
\begin{lstlisting}
int up[200001][20];
vector<vector<int>> tree(200001);
void binary_lifting(int src, int par) {
  up[src][0] = par;
  for (int i = 1; i <= 19; i++) {
    if (up[src][i - 1] != -1)
      up[src][i] = up[up[src][i-1]][i-1];
    else
      up[src][i] = -1;
  }
  for (int i = 0; i < tree[src].size(); i++)
    if (tree[src][i] != par)
      binary_lifting(tree[src][i], src);
}
int lift_node(int node, int k) {
  if (node == -1 || k == 0)
    return node;
  for (int i = 19; i >= 0; i--)
    if (k >= (1 << i))
      return lift_node(up[node][i], k-(1<<i));
}
\end{lstlisting}

\textbf{Tree Flattening (Euler Tour)}
\begin{lstlisting}
void eularDFS(ll node, ll par, ll arr[], vector<ll>&
flatTree, vector<ll> tree[], pair<ll, ll>
subTreeRange[], ll& timer)
{
  subTreeRange[node].first = ++timer;
  flatTree.push_back(arr[node]);
  for (auto& child : tree[node])
    if (child != par)
      eularDFS(child, node, arr, flatTree, tree,
      subTreeRange, timer);
  subTreeRange[node].second = timer;
}
void solve() {
  vector<ll> flatTree;
  pair<ll, ll> subTreeRange[n + 1];
  ll timer = -1;
  eularDFS(1, 0, arr, flatTree, tree,
  subTreeRange, timer);
}
\end{lstlisting}

\textbf{Heavy Light Decomposition (HLD on Tree)}
// u to v range query
\begin{lstlisting}
struct HLD {
  int n, cur_pos = 0;
  vector<int> arr, parent, size, depth, heavy,
  head, pos;
  vector<vector<int>> adj;
  FenwickTree segtree;
  HLD(int n) : n(n), arr(n + 1), parent(n + 1),
  size(n + 1), depth(n + 1), heavy(n + 1, 0), head(n +
  1) , pos(n + 1) , adj(n + 1),
  segtree(FenwickTree(n + 1)){ }
  void add_edge(int u, int v) {
    adj[u].push_back(v), adj[v].push_back(u);
  }
  void dfs(int u, int p) {
    parent[u] = p;
    size[u] = 1;
    int max_sz = 0;
    for (int v : adj[u]) {
      if (v == p)
        continue;
      depth[v] = depth[u] + 1;
      dfs(v, u);
      size[u] += size[v];
      if (size[v] > max_sz) {
        max_sz = size[v], heavy[u] = v;
      }
    }
  }
  void decompose(int u, int h) {
    head[u] = h;
    pos[u] = cur_pos++;
    segtree.add(pos[u], arr[u]); // Update
    // Fenwick tree with the +value of arr[u]
    if (heavy[u])
      decompose(heavy[u], h);
    for (int v : adj[u])
      if (v != parent[u] && v != heavy[u])
        decompose(v, v);
  }
\end{lstlisting}

\switchcolumn

\begin{lstlisting}
  int query(int u, int v) { // Range Sum Query
    int res = 0;
    while (head[u] != head[v]) {
      if (depth[head[u]] < depth[head[v]])
        swap(u, v);
      res += segtree.sum(pos[head[u]], pos[u]);
      u = parent[head[u]];
    }
    if (depth[u] > depth[v])
      swap(u, v);
    res += segtree.sum(pos[u], pos[v]);
    return res;
  }
  void addVal(int idx, int delta) {
    segtree.add(pos[idx], delta);
    arr[idx] += delta;
  }
  void editVal(int idx, int value) {
    segtree.add(pos[idx], value - arr[idx]);
    arr[idx] = value;
  }
};
\end{lstlisting}

\textbf{Centroid Decomposition}
\begin{lstlisting}
vector<vector<int>> adj;
vector<bool> is_removed;
vector<int> subtree_size;
int get_subtree_size(int node, int parent = -1) {
  subtree_size[node] = 1;
  for (int child : adj[node]) {
    if (child == parent || is_removed[child]) {
      continue; }
    subtree_size[node] +=
    get_subtree_size(child, node);
  }
  return subtree_size[node];
}
int get_centroid(int node, int tree_size, int parent
= -1) {
  for (int child : adj[node]) {
    if (child == parent || is_removed[child]) {
      continue; }
    if (subtree_size[child] * 2 > tree_size) {
      return get_centroid(child, tree_size, node);
    }
  }
  return node;
}
void build_centroid_decomp(int node = 0) {
  int centroid = get_centroid(node,
  get_subtree_size(node));
  // do something
  is_removed[centroid] = true;
  for (int child : adj[centroid]) {
    if (is_removed[child]) { continue; }
    build_centroid_decomp(child);
  }
}
\end{lstlisting}

\textit{// Ahsan's}
\begin{lstlisting}
int tot = 0;
vector<int> sz(n, 1), removed(n, 0), cenpar(n, -1);
auto dfs_sz = [&](auto &&self, int node, int pr) ->
void {
  tot += 1;
  sz[node] = 1;
  for (auto &son : g[node]) {
    if (son == pr || removed[son]) continue;
    self(self, son, node);
    sz[node] += sz[son];
  }
};
auto find_cen = [&](auto &&self, int node, int pr)
-> int {
  for (auto &son : g[node]) {
    if (son == pr || removed[son]) continue;
    if (sz[son] > tot / 2) return self(self, son, node);
  }
  return node;
};
auto decompose = [&](auto &&self, int node, int
pr) -> void {
  tot = 0;
  dfs_sz(dfs_sz, node, pr);
  int cen = find_cen(find_cen, node, pr);
  removed[cen] = 1;
  cenpar[cen] = pr;
  for (auto &son : g[cen]) {
    if (son == pr || removed[son]) continue;
    self(self, son, cen);
  }
};
decompose(decompose, 0, -1);
\end{lstlisting}

\switchcolumn

\textbf{Tree}
\begin{lstlisting}
vector<vector<int>>t,parents;
vector<int>level;
int mx;
void dfs(int node, int parent, int cLevel){
  level[node] = cLevel;
  parents[0][node] = parent;
  for(auto &x: t[node]){
    if(x!=parent) dfs(x,node, cLevel+1);
  }
}
int kthParent(int a, int k){
  for(int i = 0; i<mx; i++){
    if(a==-1) return a;
    if(k&(1<<i)) a = parents[i][a];
  }
  return a;
}
int lca(int a, int b){
  if(level[a]>level[b]){
    a = kthParent(a, level[a] - level[b]);
  }else{
    b = kthParent(b, level[b] - level[a]);
  }
  if(a==b) return a;
  for(int i = mx-1; i>=0; i--){
    if(parents[i][a]!=parents[i][b]){
      a = parents[i][a];
      b = parents[i][b];
    }
  }
  return parents[0][a];
}
int distance(int a,int b){
  return level[a]+level[b] - 2*level[lca(a,b)];
}
int main(){
  int n;
  cin>>n;
  mx = log2(n)+1;
  //resize and input
  dfs(1,-1,0);
  for(int i = 1; i<mx; i++){
    for(int j = 1; j<=n; j++){
      int prevParent = parents[i-1][j];
      if(prevParent!=-1) parents[i][j] =
      parents[i-1][prevParent];
    }
  }
}
\end{lstlisting}

\textbf{String Double Hashing}
\begin{lstlisting}
ll mod1 = 127657753, mod2 = 987654319, p1 = 31,
p2 = 37;
ll invp1 = ModPow(p1, mod1 - 2, mod1), invp2 =
ModPow(p2, mod2 - 2, mod2);
vector<pair<ll, ll>> pPow(1, {1, 1}), iPow(1, {1,
1});
void genP(ll n) {
  while ((ll)pPow.size() <= n) {
    pPow.push_back({pPow.back().xx * p1 %
    mod1, pPow.back().yy * p2 % mod2});
    iPow.push_back({iPow.back().xx * invp1 %
    mod1, iPow.back().yy * invp2 % mod2});
  }
}
struct Hash {
  int n = 0;
  vector<pair<ll, ll>> sHash;
  void genHash(string& s) {
    sHash.resize(n + 1, {0, 0}); // 0-indexed
    for (int i = 0; i < n; i++) {
      sHash[i + 1].xx = (sHash[i].xx + (s[i] - 'a' +
      1) * pPow[i].xx) % mod1;
      sHash[i + 1].yy = (sHash[i].yy + (s[i] - 'a' +
      1) * pPow[i].yy) % mod2;
    }
  }
  pair<ll, ll> getHash(int l, int r) {
    ll h1 = (sHash[r + 1].xx - sHash[l].xx + mod1)
    % mod1 * iPow[l].xx % mod1;
    ll h2 = (sHash[r + 1].yy - sHash[l].yy + mod2)
    % mod2 * iPow[l].yy % mod2;
    return {h1, h2}; // 0-indexed
  }
  pair<ll, ll> getHash() {
    return sHash[n]; // Get the hash of the whole
    // string
  }
  Hash(string& s) {
    n = s.size();
    genP(n), genHash(s);
  }
};
\end{lstlisting}

\switchcolumn
